\section{Vorlage und Spiel starten}

\begin{itemize}
    \item Kopier den ganzen Ordner SpielProjekt vom Ressourcen-Laufwerk in euren Ordner H:\\  \emph{R:/gy0187/klassen/gy0187-klasse-9b/Vorlagen/SpielProjekt}
    \item Öffne die Datei \emph{package.bluej} (manchmal heißt sie auch nur \emph{package}) in deiner Kopie der Vorlage.
    \item Mit einem Rechtsklick auf \emph{Spieler}$\rightarrow$Blockly öffnest du den Blockly-Editor. \textbf{Öffne keinesfalls Spielfeld mit Blockly!}
    \item Gehe anschließend zurück zum Hauptfenster und drücke \emph{Übersetzen} bzw. \emph{Compile}.
    \item War das erfolgreich, kannst du ein Spielfeld-Objekt und ein Spieler-Objekt erstellen, das Spielerobjekt zum Spielfeld hinzufügen und das Spiel starten (siehe Bilder unten). Wenn du möchtest, kannst du mit der Methode \emph{Spield.muenzenErzeugen(anzahl)} Münzen zum Einsammeln auf dem Spielfeld hizufügen.
\end{itemize}

\includegraphics[width=0.75\textwidth]{Aufgaben/BlocksAndJava/img/AnleitungSpielstarten.png}